\section{Limitations and Validation Boundary}
\label{sec:s2-limitations}

The MAC is an abstract discrete-event model, not an implementation of IEEE
802.11, IEEE 802.15.4, 5G NR, or another named stack. Frame sizes, PHY rate,
transmit power, carrier-sense graph, and Gilbert--Elliott parameters are
declared simulation inputs rather than radio measurements. Independent link
chains omit spatially correlated fading, capture effects, rate adaptation,
clock drift, hardware queues, and firmware latency. Offered airtime and the
energy proxy are not measured energy.

The validated selector covers one fully sensed p-persistent CSMA family and
one slotted-ALOHA cell. EXP16 confirms a formal route-by-configured-MAC
interaction only in the matched N5 Moderate cell. EXP17 supports only 1/8
core operating contexts and rejects configured MAC type as a general selector.
It is not evidence that every carrier-sensing protocol
should suppress standalone feedback or that every collision-based protocol
should enable it. Background load, hidden terminals, reverse asymmetry, and
estimator noise are reported boundaries, not a complete operational design
space. The $N=20$ causal ACK effect and adaptive--piggyback equivalence are not
identified.

The primary $N=5$ experiment uses 6-DOF follower dynamics, but larger swarms
and all stated formation theorems use the double-integrator subsystem. The ISS
result assumes inactive saturation and an undirected grounded graph; the
finite-horizon safety certificate is conditional and presently conservative.
Raw zero failure counts are observed guards, not population guarantees.

ACK integrity assumes honest cumulative feedback. Sequence checks, rollback
rejection, and future-generation rejection detect protocol violations under
loss and delay, but no spoofing, replay attacker, jamming adversary, or key
management is modeled. Security claims require an authenticated feedback
design and separate threat-driven experiments.

Finally, the study has no trace, radio, HIL, or flight validation. The frozen
EXP17 continuation gate fails, so hardware procurement for the current general
selector is deferred rather than presented as the automatic next step. The
measured absolute-time trace and two-node airtime/energy interfaces remain
available for a future redesigned policy or a deliberately narrow base-cell
mechanism study.

\section{Conclusion}
\label{sec:s2-conclusion}

Receiver-confirmed semantic triggering remains causal on a finite shared
medium, but causality does not make every acknowledgement useful. Broadcast
aggregation, scaled access, cumulative piggybacking, and standalone feedback
interact through future DATA demand and contention. The initial negative
holdout, factorial repair, focal-decision replay, disjoint validation, and a
matched factorial interaction expose one bounded mechanism: piggyback is
favored in the tested base CSMA cell while a standalone path is favored in its
matched ALOHA cell.

The broader operating-envelope holdout supports only 1/8 contexts and rejects
configured MAC type as a general selector. Together with the $N=20$ support
failure and adverse focal effects, this prevents a universal ACK-value or
deployment claim. The contribution is instead a causal account of feedback
mediation and an experimentally identified boundary showing why confirmation
lead, MAC category, and communication cost cannot individually determine
closed-loop value.
